% Blueprint for Minimum Balanced Bipartitions of Planar Triangulations
% Generated from minbal_tex/ex_article.tex

\section{Introduction}

\begin{definition}[Balanced Bipartition]
A balanced bipartition of a graph $G$ is a bipartition $(V_1, V_2)$ of $V(G)$ where $V_1$ and $V_2$ differ in size by at most 1.
\end{definition}

\begin{definition}[Minimum Balanced Bipartition]
A minimum balanced bipartition of $G$ is a balanced bipartition $(V_1, V_2)$ of $V(G)$ with the minimum number $e(V_1, V_2)$ of edges with ends in both $V_1$ and $V_2$. 
\end{definition}

\begin{theorem}[Main Theorem]
\label{thm:main}
If $G$ is a plane triangulation, then there exists a balanced bipartition $(V_1, V_2)$ of $V(G)$ such that both $G[V_1]$ and $G[V_2]$ are connected near-triangulations, and the total number of blocks in $G[V_1]$ and $G[V_2]$ exceeds the total number of internal vertices by at most 2.
\end{theorem}

\section{Preliminaries}

\begin{definition}[Near-Triangulation]
A near-triangulation is a plane graph $G$ such that every face of $G$ except at most one is bounded by three edges.
\end{definition}

\begin{definition}[Internal Vertex]
Any vertex of $G$ not lying on the boundary of the infinite face of $G$ is said to be an internal vertex of $G$.
\end{definition}

\begin{proposition}[Edge Count in Near-Triangulation]
\label{prop:edges_cpnt}
If $G$ is a connected near-triangulation with $|V(G)| \geq 2$, $b$ blocks, and $i$ internal vertices, then $|E(G)| = 2|V(G)|-2+i-b$.
\end{proposition}

\begin{corollary}[Edge Cut in Bipartition]
\label{cor:cut_cpnt}
If $G$ is a plane triangulation with $|V(G)| \geq 4$, and $(V_1, V_2)$ is a bipartition of $V(G)$ with $|V_1|, |V_2| \geq 2$ such that both $G[V_1], G[V_2]$ are near-triangulations together containing a total of $b$ blocks and $i$ internal vertices, then $e(V_1, V_2) = |V(G)| + b - i - 2$.
\end{corollary}

\section{Partitioning Lemmas}

\begin{lemma}[4-Connected Partitioning]
\label{lem:4con}
If $G$ is a 4-connected plane triangulation or $K_4$, and $e_1, e_2$ are disjoint edges in $G$, then for every $2 \leq k \leq |V(G)|-2$ there exists a bipartition $(V_1, V_2)$ of $V(G)$ such that:
\begin{enumerate}
\item $|V_1| = k, |V_2| = |V(G)|-k$,
\item $e_1 \in E(G[V_1]), e_2 \in E(G[V_2])$, and
\item both $G[V_1], G[V_2]$ are biconnected near-triangulations.
\end{enumerate}
\end{lemma}

\begin{corollary}[Conjecture for 4-Connected Graphs]
\label{cor:conj_4con}
If $G$ is a 4-connected planar graph or $G \subseteq K_4$, then for every $1 \leq k \leq |V(G)|-1$ there exists a bipartition $(V_1, V_2)$ of $V(G)$ with $|V_1| = k$ such that $e(V_1, V_2) \leq |V(G)|$.
\end{corollary}

\begin{lemma}[3-Connected Partitioning]
\label{lem:3con1}
If $G$ is a plane triangulation and the vertex-set ${a, b, c} \subseteq V(G)$ bounds the infinite face of $G$, then for every $2 \leq k \leq |V(G)|-1$ there exists a bipartition $(V_1, V_2)$ of $V(G)$ such that:
\begin{enumerate}
\item $|V_1| = k, |V_2| = |V(G)|-k$,
\item ${a, b} \subseteq V_1, c \in V_2$,
\item $G[V_1]$ is a biconnected near-triangulation with the edge $ab$ on the boundary of the infinite face, and
\item $G[V_2]$ is a connected near-triangulation with the vertex $c$ on the boundary of the infinite face.
\end{enumerate}
\end{lemma}

\begin{lemma}[Partition Refinement]
\label{lem:3con2}
Let $G$ be a plane triangulation such that $|V(G)| \geq 4$ and the vertex-set ${a, b, c} \subset V(G)$ bounds the infinite face of $G$. If $(V_1, V_2)$ is a bipartition of $V(G)$ such that ${a, b} \subseteq V_1, c \in V_2$, $|V_2| \geq 2$, and both $G[V_1]$ and $G[V_2]$ are connected near-triangulations, then there exists a vertex $v \in V_2-c$ such that both $G[V_1 \cup v]$ and $G[V_2-v]$ are connected near-triangulations, and the total number of blocks in $G[V_1 \cup v]$ and $G[V_2-v]$ exceeds that in $G[V_1]$ and $G[V_2]$ by at most 1.
\end{lemma}

\section{Setup and Construction}

\begin{proposition}[Warmup Case]
\label{prop:warmup}
If $G$ is a plane triangulation with $|V(G)|=n$ and $T \subseteq G$ is a separating triangle such that each of the two components of $G-V(T)$ contains at least $\lfloor \frac{n}{2} \rfloor-1$ vertices, then there exists a balanced bipartition $(V_1, V_2)$ of $V(G)$ such that $G[V_1]$ is a biconnected near-triangulation and $G[V_2]$ is a connected near-triangulation with the number of blocks exceeding the total number of internal vertices in $G[V_1]$ and $G[V_2]$ by at most 1.
\end{proposition}

\begin{proposition}[Sink Degree Case]
\label{prop:onedegsink}
If $G$ is a plane triangulation with $|V(G)|=n$ such that the sink in its directed std-tree has degree 1, then there exists a balanced bipartition $(V_1, V_2)$ of $V(G)$ such that both $V_1$ and $V_2$ induce biconnected near-triangulations.
\end{proposition}

\begin{proposition}[Sink with 4 Vertices]
\label{prop:sonly4}
If $G$ is a plane triangulation with $|V(G)|=n$ such that the sink triangulation $S$ of $G$ contains exactly 4 vertices, then there exists a balanced bipartition $(V_1, V_2)$ of $V(G)$ such that both $V_1$ and $V_2$ induce connected near-triangulations, and the total number of blocks in $G[V_1]$ and $G[V_2]$ exceeds the total number of internal vertices by at most 2.
\end{proposition}

\section{Algorithm and Proof Structure}

The proof follows a constructive approach using the following key components:

\begin{enumerate}
\item {\bf Separating Triangle Decomposition}: Decompose the plane triangulation into 4-connected pieces using separating triangles.
\item {\bf Navigation Lemmas}: Use Lemmas \ref{lem:4con} and \ref{lem:3con1} to navigate the sink triangulation and build the bipartition incrementally.
\item {\bf Partition Refinement}: Use Lemma \ref{lem:3con2} to refine partitions while maintaining connectivity and near-triangulation properties.
\item {\bf Special Case Handling}: Handle special cases where the sink triangulation has degree 1 or exactly 4 vertices.
\end{enumerate}

The algorithm {\bf Construct Bipartition} implements this approach:

\begin{enumerate}
\item Initialize with triangles having largest smaller sides
\item Navigate the sink triangulation using partitioning lemmas
\item Handle split configurations where direct partitioning is needed
\item Ensure the final bipartition satisfies the block count constraint
\end{enumerate}

\section{Conclusion}

The main theorem \ref{thm:main} is proven by:
\begin{itemize}
\item Reducing to the case where no separating triangle splits the graph evenly
\item Using the separating triangle decomposition to identify a 4-connected sink
\item Applying the partitioning lemmas to construct the desired balanced bipartition
\item Handling special cases through dedicated propositions
\item Verifying that the constructed bipartition satisfies all required properties
\end{itemize}

This confirms the folklore conjecture that for any planar graph $G$, a minimum balanced bipartition $(V_1, V_2)$ of $V(G)$ has $e(V_1, V_2) \leq |V(G)|$.